Consider a coordinate system with the two given points lying on the horizontal axis
at coordinates $A: \paren{-a, 0}$ and $B: \paren{a, 0}$, for some $a \neq 0$.
Let $C: \paren{b, c}$ be the center of the given circle,
and let $s > 0$ be the radius. \newline
The wanted circle passes through $A$ and $B$,
so its center must lie on the perpendicular bisector of the two points,
which is the vertical axis;
so it has coordinates $P: \paren{0, d}$,
for some $d \in \mathbb{R}$.
Let $r$ be its radius.
The tangency condition between the two circles is equivalent to
$\dist{P}{C} = r + s$ (for external tangency)
or $\dist{P}{C} = r - s$ (when the first circle is internally tangent to the wanted circle)
or $\dist{P}{C} = s - r$ (when the wanted circle is internally tangent to the first circle).
All in all, the tangency conditions can be summarised as
$\dist{P}{C}^2 = \paren{r \pm s}^2$.
Together with the condition $r = \dist{P}{A} = \dist{P}{B}$,
this yields the following system of equations,
in the unknowns $d$ and $r$:
\begin{align*}
\begin{cases}
 r^2 = a^2 + d^2 \\
 r^2 + s^2 \pm 2 r s = b^2 + \paren{c - d}^2
\end{cases}
\end{align*}
Substituting $a^2 + d^2$ for $r^2$ in the second equation,
and canceling the term $d^2$ yields:
\begin{equation*}
 a^2 + s^2 \pm 2 r s = b^2 + c^2 - 2 c d
\end{equation*}
So:
\begin{equation*}
 \mp r = \frac{c}{s} d + \frac{1}{2 s} \paren{a^2 - b^2 - c^2 + s^2}
\end{equation*}
Substituting in the first equation yields a quadratic equation in $d$:
\begin{equation*}
 \paren{c^2 - s^2} d^2
 + c \paren{a^2 - b^2 - c^2 + s^2} d
 + \frac{1}{4} \paren{a^2 - b^2 - c^2 + s^2}^2 - a^2 s^2
 = 0
\end{equation*}

If $c^2 = s^2$, the equation is actually linear:
\begin{equation*}
 c \paren{a^2 - b^2} d
 + \frac{1}{4} \paren{a^2 - b^2}^2 - a^2 s^2
 = 0
\end{equation*}
Note that $a^2 - b^2$ cannot be $0$, since it would imply $a = 0$ or $s = 0$.
Also, $c \neq 0$, since $c^2 = s^2 \neq 0$.
The solution in this case is then:
\begin{equation*}
d = \frac{a^2 s^2 - \mfrac{1}{4} \paren{a^2 - b^2}^2}{c \paren{a^2 - b^2}}
   = \frac{a^2 c}{a^2 - b^2} - \frac{a^2 - b^2}{4 c}
\end{equation*}

Observe that $c^2 = s^2$ if and only if the line $AB$ is tangent to the first circle.
In this case, the line $AB$ can be considered a degenerate second solution circle.

\vspace{3ex}
In the case $c^2 \neq s^2$, the solutions are:
\begin{align*}
 d &= \frac{- \, c \, \paren{a^2 - b^2 - c^2 + s^2}
     \pm \sqrt{c^2 \paren{a^2 - b^2 - c^2 + s^2}^2
               - \paren{c^2 - s^2} \paren{a^2 - b^2 - c^2 + s^2}^2
               + 4 \paren{c^2 - s^2} a^2 s^2}}
     {2 \paren{c^2 - s^2}}
   \\[0.3cm]
   &= \frac{- \, c \, \paren{a^2 - b^2 - c^2 + s^2}
     \pm \, s \, \sqrt{\paren{a^2 - b^2 - c^2 + s^2}^2
                       + 4 a^2 c^2 - 4 a^2 s^2}}
     {2 \paren{c^2 - s^2}} =
   \\[0.3cm]
   &= \frac{- \, c \, \paren{a^2 - b^2 - c^2 + s^2}
     \pm \, s \, \sqrt{\paren{a^2 + b^2 + c^2 - s^2}^2
                       - 4 a^2 b^2}}
     {2 \paren{c^2 - s^2}} =
   \\[0.3cm]
   &= \frac{- \, c \, \paren{a^2 - b^2 - c^2 + s^2}
     \pm \, s \, \sqrt{\paren{\paren{a + b}^2 + c^2 - s^2}
                       \paren{\paren{a - b}^2 + c^2 - s^2}}}
     {2 \paren{c^2 - s^2}} =
   \\[0.3cm]
   &= \frac{- \, c \, \paren{a^2 - b^2 - c^2 + s^2}
     \pm \, s \, \sqrt{\paren{\dist{A}{C}^2 - s^2}
                       \paren{\dist{B}{C}^2 - s^2}}}
     {2 \paren{c^2 - s^2}}
\end{align*}
Observe that
$\paren{\paren{a + b}^2 + c^2 - s^2} \paren{\paren{a - b}^2 + c^2 - s^2}$
is positive if and only if $A$ and $B$ are both inside the first circle or both outside.
It is $0$ if and only if at least one among the points $A$ and $B$ lies on the circle.

In the general case, there are two circles that satisfy the stated conditions,
and they are each uniquely determined by the coordinate $d$.

If the problem asks to perform the construction with only straightedge and compass,
then it can be achieved by solving the quadratic equation for $d$ geometrically,
via, for example, the Carlyle circle construction,
since all coefficients are constructible.
Alternatively, the construction can be carried out by following the formula for $d$,
since it only involves operations that can be performed with straightedge and compass.
Note that it can be assumed that, for example, $s = 1$, since everything is scale-invariant.

In the case $s^2 = c^2$, the construction can also be performed with straightedge and compass,
since the formula for $d$ involves constructible operations only.
